\mathquote{Mladý muži, v matematice věcem nerozumíte. Jen si na ně zvyknete.}{John von Neumann}

Maturitu máte za sebou, nebo před sebou, ale to teď nevadí. V této části knihy se podíváme, kam matematika vede dál. Nic z toho, co tu najdete, u maturity potřebovat nebudete. Tak proč to číst?

\section*{Proč se do toho pustit už teď}

Pokud se chystáte studovat matematiku, fyziku, informatiku, nějaký technický obor, ekonomii nebo třeba chemii, tyto věci vás \textbf{stejně čekají}. Lineární algebra a matematická analýza (derivace, integrály, řady a funkce více proměnných) jsou na většině vysokých škol tohoto typu hned v prvním semestru. Diferenciální rovnice, komplexní analýza a teorie pravděpodobnosti obvykle přijdou o rok nebo dva později. A patří mezi předměty, ze kterých mají studenti prvního ročníku největší respekt.

Na vysoké škole se tyto předměty učí rychle a formálně: definice, věta, důkaz, a další týden nová kapitola. Na obrázky a intuici obvykle nezbývá čas. Kdo ale ty pojmy už jednou viděl, třeba jen zběžně a s obrázky, má obrovskou výhodu. Na přednášce pak neluští, \textit{co} se děje, ale může sledovat, \textit{proč} to platí. Tahle část knihy se snaží být přesně tím prvním setkáním: stejně názorná jako maturitní část, jen o patro výš.

A pokud nic takového studovat nechcete, můžete tu najít odpovědi na otázky, které vás možná napadly už při maturitě. Proč má kalkulačka v sobě sinus, když umí jen sčítat a násobit? Proč $\mathrm{e}^{i\pi} = -1$? Proč se nedá vyřešit rovnice pátého stupně? Jak se počítá plocha pod Gaussovou křivkou, když k ní neexistuje primitivní funkce?

\section*{Jak tuto část číst}

\begin{itemize}
  \item Kapitoly stojí na první části knihy. Kde se něco používá, je tam odkaz, takže se můžete kdykoli vrátit.
  \item Doporučené pořadí je lineární algebra, řady, funkce více proměnných a diferenciální rovnice, protože na sebe navazují. Teorii čísel můžete číst kdykoli. Komplexní analýza, teorie míry a fyzika jsou nejnáročnější a hodí se nakonec.
  \item Nemusíte rozumět všemu napoprvé. Důkazy klidně při prvním čtení přeskočte a vraťte se k nim později. Matematika se na vysoké škole nechápe za jeden večer, ale postupně, opakovaným vracením.
  \item Cvičení na konci kapitol mají řešení. Zkuste je, i když se vám nebudou dařit. Právě u těch, které nejdou, se člověk naučí nejvíc.
\end{itemize}

Citát na začátku je od jednoho z nejchytřejších lidí 20. století. Pokud vám něco v této části bude připadat nepochopitelné, je to normální a von Neumannovi to tak zjevně připadalo taky. Po čase si zvyknete, a zvyknout si je první krok k porozumění.
